\begin{question}

  The Hamiltonian operator for a two-state system is given by
  \begin{equation*}
    H = a(\dyad 1 - \dyad 2 + \dyad{1}{2} + \dyad{2}{1}),
  \end{equation*}

  where $a$ is a number with the dimension of energy. Find the energy
  eigenvalues and the corresponding energy eigenkets (as linear
  combinations of $\ket 1$ and $\ket 2$).
\end{question}

\begin{answer}

  It's worth noting that the fact that it's a Hamiltonian is totally
  irrelevant to the problem.  It would be just as meaningful to the
  $MaraLagoFace$ operator which transforms $\hat{MLF}\ket{prettyWoman}
  = \ket{frankenFascist}$.  So long as you have a basis in which to
  express it (let's say the $\ket{brutish}$ and $\ket{bulbous}$ basis
  kets) then you can put it in matrix form and find the eigenkets
  (say, $\ket{carWreckVictim}$ for example or
  $\ket{agingBeautyQueen}$).
  
  Manually expressing the operator as a matrix, we find that:

  \begin{equation}
    \tag{1.73}
    X \dot{=}
    \begin{pmatrix}
      \mel{\alpha^{(1)}}{X}{\alpha^{(1)}} & \mel{\alpha^{(1)}}{X}{\alpha^{(2)}} & \ldots \\
      \mel{\alpha^{(2)}}{X}{\alpha^{(1)}} & \mel{\alpha^{(2)}}{X}{\alpha^{(2)}} & \ldots \\
      \vdots & \vdots & \ddots\\
    \end{pmatrix}
  \end{equation}

  \begin{equation}
    \begin{split}
      \mel{1}{H}{1}
      &= \mel{L}{H}{R}\\
      &=a(
      \braket{L}{1}\braket{1}{R}
      - \braket{L}{2}\braket{2}{R}
      + \braket{L}{1}\braket{2}{R}
      + \braket{L}{2}\braket{1}{R}
      )
      \\
      &=a(
      \braket{1}{1}\braket{1}{1}
      - \braket{1}{2}\braket{2}{1}
      + \braket{1}{1}\braket{2}{1}
      + \braket{1}{2}\braket{1}{1}
      )
      \\
      &=
      a(
      1
      - 0
      + 0
      + 0
      )
      \\
      &= a
    \end{split}
  \end{equation}

  \begin{equation}
    \begin{split}
      \mel{1}{H}{2}
      &= \mel{L}{H}{R}\\
      &=a(
      \braket{L}{1}\braket{1}{R}
      - \braket{L}{2}\braket{2}{R}
      + \braket{L}{1}\braket{2}{R}
      + \braket{L}{2}\braket{1}{R}
      )
      \\
      &=a(
      \braket{1}{1}\braket{1}{2}
      - \braket{1}{2}\braket{2}{2}
      + \braket{1}{1}\braket{2}{2}
      + \braket{1}{2}\braket{1}{2}
      )
      \\
      &=
      a(
      0
      - 0
      + 1
      + 0
      )
      \\
      &= a
    \end{split}
  \end{equation}

  \begin{equation}
    \begin{split}
      \mel{2}{H}{1}
      &= \mel{L}{H}{R}\\
      &=a(
      \braket{L}{1}\braket{1}{R}
      - \braket{L}{2}\braket{2}{R}
      + \braket{L}{1}\braket{2}{R}
      + \braket{L}{2}\braket{1}{R}
      )
      \\
      &=a(
      \braket{2}{1}\braket{1}{1}
      - \braket{2}{2}\braket{2}{1}
      + \braket{2}{1}\braket{2}{1}
      + \braket{2}{2}\braket{1}{1}
      )
      \\
      &=
      a(
      0
      - 0
      + 0
      + 1
      )
      \\
      &= a
    \end{split}
  \end{equation}

  \begin{equation}
    \begin{split}
      \mel{2}{H}{2}
      &= \mel{L}{H}{R}\\
      &=a(
      \braket{L}{1}\braket{1}{R}
      - \braket{L}{2}\braket{2}{R}
      + \braket{L}{1}\braket{2}{R}
      + \braket{L}{2}\braket{1}{R}
      )
      \\
      &=a(
      \braket{2}{1}\braket{1}{2}
      - \braket{2}{2}\braket{2}{2}
      + \braket{2}{1}\braket{2}{2}
      + \braket{2}{2}\braket{1}{2}
      )
      \\
      &=
      a(
      0
      - 1
      + 0
      + 0
      )
      \\
      &= -a
    \end{split}
  \end{equation}

  \begin{equation}
    H \dot{=}
    a\begin{pmatrix}
      1 & 1\\
      1 & -1\\
    \end{pmatrix}
  \end{equation}

  OK...if you're like me, then you haven't worked a basic Eigenvalue
  problem since college (well...I guess some of you still \emph{in}
  college) and have been leaning pretty heavily on the Eig function
  (which, incidentally, when you go source-diving seems to be
  \emph{defined} inside the header file...not just like putting a stub
  there so that the compiler knows what to expect...no...the actual
  source for it seems to live within the header file...C++ is an
  abomination... use Rust or Julia... or Python... or Gauss-help-you even
  straight C is a better choice).  No way around it, so let's dive
  right in.

  
  \begin{equation}
    \begin{split}
      0 &= det(H - \lambda I)\\
      &= (a - \lambda)(-a - \lambda) - a^2\\
      &= -a^2 - a\lambda + a\lambda + \lambda^2 - a^2\\
      &= -a^2 + \lambda^2 - a^2\\
      &= -2a^2 + \lambda^2\\
      \implies \lambda^2 &= 2a^2\\
      \implies \lambda &= \pm a \sqrt{2}\\
    \end{split}
  \end{equation}

  Now let's plug that into our characteristic equations

  \begin{equation}
    \begin{split}
      H
      \begin{bmatrix}
        x_1\\
        y_1\\
      \end{bmatrix}
      &= \lambda_1
      \begin{bmatrix}
        x_1\\
        y_1\\
      \end{bmatrix}\\
      H
      \begin{bmatrix}
        x_2\\
        y_2\\
      \end{bmatrix}
      &= \lambda_2
      \begin{bmatrix}
        x_2\\
        y_2\\
      \end{bmatrix}\\
    \end{split}
  \end{equation}

  Yeah, you can do the row reduction thing, but where's the fun in
  that?  Translating into a system of equations, and taking them one
  at a time, we find:

  \begin{equation}
    \begin{alignedat}{2}
      & & ax_1 + ay_1 &= \lambda_1x_1\\
      \implies \quad & & ax_1 + ay_1 &= a\sqrt{2}x_1\\
      \implies \quad & & x_1 + y_1 &= \sqrt{2}x_1\\
      \implies \quad & & y_1 &= (\sqrt{2} - 1)x_1\\
      & & ax_1 - ay_1 &= \lambda_1y_1\\
      \implies \quad & & ax_1 - ay_1 &= a\sqrt{2}y_1\\
      \implies \quad & & x_1 - y_1 &= \sqrt{2}y_1\\
      \implies \quad & & x_1 &= (\sqrt{2} + 1)y_1\\
      & & ax_2 + ay_2 &= \lambda_2x_2\\
      \implies \quad & & ax_2 + ay_2 &= -a\sqrt{2}x_2\\
      \implies \quad & & x_2 + y_2 &= -\sqrt{2}x_2\\
      \implies \quad & & y_2 &= -(\sqrt{2} + 1)x_2\\
      & & ax_2 - ay_2 &= \lambda_2y_2\\
      \implies \quad & & ax_2 - ay_2 &= -a\sqrt{2}y_2\\
      \implies \quad & & x_2 - y_2 &= -\sqrt{2}y_2\\
      \implies \quad & & x_2 &= (1 - \sqrt{2})y_2\\
    \end{alignedat}
  \end{equation}

  You'll note that all of the factors of $a$ go away because
  eigenvectors are relative values that get normalized later, so
  constant factors go away.  It's also pretty easy to show since the
  definition of $H$ includes the constant factor of $a$.  In general,
  it's easier to pull the factor out before we go through the effort.
  If we pick some numbers out of thin air to get the ball rolling, say
  $x_1 \equiv 1$ \& $y_2 \equiv 1$, we find that the non-normalized
  values are:
  
  \begin{equation}
    \begin{split}
      x_1 &= 1\\
      y_1 &= (\sqrt{2} - 1)x_1\\
      &= \sqrt{2} - 1\\
      y_2 &= 1\\
      x_2 &= 1 - \sqrt{2}y_2\\
      &= 1 - \sqrt{2}\\
    \end{split}
  \end{equation}

  Solving for the normalization coefficients:
  
  \begin{equation}
    \begin{split}
      N_1 &= |\vec{a_1}|\\
      &= \sqrt{\vec{a_1} \cdot \vec{a_1}}\\
      &= \sqrt{4 - 2\sqrt{2}}\\
      N_2 &= |\vec{a_2}|\\
      &= \sqrt{\vec{a_2} \cdot \vec{a_2}}\\
      &= \sqrt{4 - 2\sqrt{2}}\\
      N &= N_1 = N_2 = \sqrt{4 - 2\sqrt{2}}\\
    \end{split}
  \end{equation}

  Running a quick check with our earlier equations, we find that:

  \begin{equation}
    \begin{alignedat}{2}
      & & ax_1 + ay_1 &= \lambda_1x_1\\
      \implies \quad & & a + ay_1 &= a\sqrt{2}\\
      \implies \quad & & 1 + y_1 &= \sqrt{2}\\
      \implies \quad & & 1 + \sqrt{2} - 1 &= \sqrt{2}\\
      \implies \quad & & \sqrt{2} &= \sqrt{2}\qed\\
    \end{alignedat}
  \end{equation}

  We'll also need to restore the factor of $a$ that we dropped earlier
  (since it's just a constant multiplied by the matrix we solved for).
  Don't worry, I missed it the first few times as well.

  \begin{equation}
    \begin{split}
      \frac H a &= \sum_i \lambda_i \dyad{a_i}\\
      &\dot{=}\sum_i \lambda_i \begin{bmatrix}x_1\\y_1\\\end{bmatrix}\begin{bmatrix}x_1&y_1\\\end{bmatrix}\\
          &=
          \frac{\sqrt{2}}{N^2}
          \begin{bmatrix}1\\\sqrt{2} - 1\\\end{bmatrix}\begin{bmatrix}1&\sqrt{2} - 1\\\end{bmatrix}
              - \frac{\sqrt{2}}{N^2}
              \begin{bmatrix}1 - \sqrt{2}\\1\\\end{bmatrix}\begin{bmatrix}1 - \sqrt{2}&1\\\end{bmatrix}
    \end{split}
  \end{equation}

\end{answer}
